\documentclass[12pt]{article}
\usepackage[margin=1in]{geometry}
\usepackage{setspace}
\usepackage{url}
\usepackage{footmisc}
\usepackage{fancyhdr}  % For custom headers/footers

\doublespacing

\pagestyle{fancy}
\fancyhf{}  % Clear all header and footer fields
\fancyhead[R]{Vaught \thepage}  % Top right: Lastname PageNumber

\begin{document}

Debates over whether physicians should adopt the \textit{classical} Hippocratic Oath, translated by Ludwig Edelstein in 1943, or the \textit{modern} version, penned by Louis Lasagna in 1964, have persisted for decades.\footnote{See Edelstein; Lasagna.} Numerous commentators have presumed that explicit nods to patient autonomy, empathy, and cultural sensitivity in the modern oath naturally render the ancient text obsolete. Yet, evidence from both historical inquiry and contemporary ethical failures—including high-profile abuses of research subjects—indicates that revised wording alone does not guarantee ethical practice. Indeed, some scholars have argued that the classical oath, despite references to Greek deities and perceived prohibitions on abortion, euthanasia, or even surgery, may hold special relevance for modern times, precisely because of its unwavering moral clarity and its insistence on “do no harm.” The classical text—far from being a mere historical curiosity—continues to underscore a physician’s duty to address not only individual patient welfare but also broader injustices in health care.\footnote{See Miles 99--105.}

\bigskip

Some have alleged that swearing by Greek gods disqualifies the classical oath from use in today’s religiously pluralistic world, suggesting that lines invoking “Apollo the Physician and Asclepius and Hygieia and Panaceia” introduce exclusivist or even pagan elements.\footnote{See Edelstein 1--2.} Nevertheless, theologians and historians have shown that this invocation actually frames the physician’s duties within a transcendent moral order, highlighting dedication to healing and humility before higher principles.\footnote{See Jacquart and Thomasset 18--25.} Many versions of the classical oath have already replaced the ancient deities with monotheistic references—or even secular affirmations—without betraying its intent to remind practitioners that medical power must be tempered by ethical limits.\footnote{See Veatch and Mason 86--105.} In fact, rather than enforcing any single theology, the classical oath’s opening serves to emphasize a divine or higher accountability, which can include Christian, Jewish, or Islamic conceptions of God—or even broader ethical commitments—if one so chooses.\footnote{Ibid.} Thus, instead of rendering the oath obsolete, these theistic references underscore a universal message that physicians stand answerable to something beyond personal interests, whether that be conscience, the divine, or societal well-being.

\bigskip

Others have protested that the classical oath categorically forbids core procedures such as abortion and surgery, thereby clashing with contemporary medical standards.\footnote{See Omran; Veatch and Mason 90--94.} Critics point to phrases like “I will not give a drug that is deadly to anyone if asked,” interpreting it as a strict ban on euthanasia, and “I will not give a woman a destructive pessary,” reading it as an unequivocal condemnation of abortion.\footnote{Ibid.} Yet, scholars have clarified that ancient Greek physicians often feared being co-opted as political assassins or “physician poisoners,” so the oath’s statements targeted malicious uses of medical expertise—rather than thoughtful end-of-life care.\footnote{See Miles 42--45; Porter 54--59.} Similarly, the text does not ban all surgery: “I will not use the knife, not even on sufferers from stone, but will withdraw in favor of such men as are practitioners of this work,” likely envisioned the emerging concept of specialist surgeons, acknowledging that some forms of invasive intervention lay beyond general practice.\footnote{See Edelstein 2.} Modern specialists—such as oncologists, radiologists, and cardiothoracic surgeons—exemplify the oath’s foresight in urging physicians to work collaboratively within their specific domains of expertise.\footnote{See Siraisi 78--80.} Accordingly, the classical oath’s prohibitions, properly contextualized, champion patient safety rather than stifling scientific progress.

\bigskip

A further critique holds that the classical oath focuses too narrowly on harm prevention while neglecting empathy, social justice, and informed consent—qualities that the modern oath explicitly affirms.\footnote{See Lasagna.} Indeed, Lasagna’s version proclaims, “I will remember that there is art to medicine as well as science,” and includes language emphasizing patient privacy and respect for autonomy.\footnote{Ibid.} However, the classical oath’s requirement that physicians not exploit vulnerable patients—“especially from sexual acts,” as one passage states—directly enshrines a duty to protect patient dignity and welfare.\footnote{See Edelstein 3.} Moreover, its insistence on confidentiality—“Whatever I may see or hear in the course of treatment… I will keep to myself”—laid the groundwork for modern privacy principles, well before legal frameworks codified them.\footnote{Ibid.} Recent scholarship has also noted that “do no harm” expands naturally into challenging injustices both inside and outside the clinic: if harm includes lack of access to care or prohibitive costs of drugs, then the oath demands advocacy for more equitable health systems.\footnote{See Miles 110--115; Porter 66--69.} In this manner, the classical oath, far from being ethically rudimentary, can drive a profound commitment to justice, empathy, and patient-centered reform.

\bigskip

In conclusion, while the \textit{modern} Hippocratic Oath offers praiseworthy clarity on certain contemporary ethical concerns, the \textit{classical} oath remains a potent and adaptable code that may be more relevant now than ever. Its succinct and timeless injunction—“do no harm”—has proven resilient amid dramatic advances in surgery, pharmacology, and medical technology, and it underscores the physician’s obligation to resist harmful pressures from political, financial, or ideological forces.\footnote{See Jacquart and Thomasset 122--125.} Where the modern oath enumerates specific commitments, the classical text compels continual ethical reflection, urging physicians to confront new “deities” such as profit motives, insurance complexities, or institutional biases that risk undermining patient welfare. Ultimately, adopting a slightly updated classical oath—shorn of archaic language but preserving its moral core—may better emphasize that medicine’s fundamental mission transcends both historical boundaries and immediate policy debates, uniting physicians under a shared promise to protect patients and society from harm and injustice.

\bigskip

\begin{thebibliography}{9}

\bibitem{Edelstein}
Ludwig Edelstein, \textit{The Hippocratic Oath: Text, Translation, and Interpretation} (Baltimore: Johns Hopkins University Press, 1943), accessed via PBS website, \url{https://www.pbs.org/wgbh/nova/article/hippocratic-oath-today/}.

\bibitem{Lasagna}
Louis Lasagna, “Hippocratic Oath,” \textit{Tufts University School of Medicine} (1964), accessed via PBS website, \url{https://www.pbs.org/wgbh/nova/article/hippocratic-oath-today/}.

\bibitem{Miles}
Stephen H. Miles, \textit{The Hippocratic Oath and the Ethics of Medicine} (Oxford: Oxford University Press, 2005).

\bibitem{Jacquart}
Danielle Jacquart and Claude Thomasset, \textit{Sexuality and Medicine in the Middle Ages} (Princeton, NJ: Princeton University Press, 1988).

\bibitem{Veatch}
Robert M. Veatch and Carol Mason, “Hippocratic vs. Judeo-Christian Medical Ethics,” \textit{Journal of Religious Ethics} 15, no. 1 (1987): 86–105.

\bibitem{Omran}
Adel Omran, “Family Planning in the Legacy of Islam,” in \textit{Routledge Library Editions: Women, Family and Kinship} 9 (1992): 23–24.

\bibitem{Porter}
Roy Porter, \textit{The Greatest Benefit to Mankind: A Medical History of Humanity from Antiquity to the Present} (New York: W. W. Norton, 1999).

\bibitem{Siraisi}
Nancy G. Siraisi, \textit{Medieval \& Early Renaissance Medicine: An Introduction to Knowledge and Practice} (Chicago: University of Chicago Press, 1990).

\end{thebibliography}

\end{document}